\documentclass[11pt]{article}
\usepackage[margin=2.6cm]{geometry}
\usepackage{amsmath,amssymb,amsthm}
\usepackage{enumitem}
\usepackage{xr}
\externaldocument[M-]{model}
\externaldocument[L-]{localisation}
\setlist{itemsep=1pt,topsep=3pt}

\newcommand{\C}{\mathcal{C}}
\newcommand{\Path}{\mathrm{Path}}
\newcommand{\U}{\mathsf{U}}
\newcommand{\El}{\mathsf{El}}
\newcommand{\id}{\mathrm{id}}
\newcommand{\D}{\mathsf{D}}
\newcommand{\cosk}{\mathrm{cosk}}
\newcommand{\Lam}{\Lambda}

\theoremstyle{plain}
\newtheorem{theorem}{Theorem}[section]
\newtheorem{lemma}[theorem]{Lemma}
\newtheorem{proposition}[theorem]{Proposition}
\newtheorem{corollary}[theorem]{Corollary}
\theoremstyle{definition}
\newtheorem{definition}[theorem]{Definition}
\newtheorem{example}[theorem]{Example}
\theoremstyle{remark}
\newtheorem{remark}[theorem]{Remark}

\title{Finite approximations of the localisation\\[3pt]
\large a complement to \emph{The localisation of the poset model}}
\author{}
\date{21 October 2026}

\begin{document}
\maketitle

\begin{abstract}
\noindent
For a type $A$ of the poset model, $\D_NA$ is defined as $\D A$
(\texttt{localisation.tex}) with the families indexed only by the maps
$[k]\to X$ with $k\le N$. We show that $\D_N A=\D(K_NA)$, where $K_NA$ is
the $N$-coskeleton of $A$: a type over the posets whose restriction to the
linear posets is the simplicial coskeleton of $i^*A$. So $\D_NA$ is fibrant
when $A$ is, by Theorem~\ref{L-thm:DA} of \texttt{localisation.tex}
applied to $K_NA$. Spheres of dimension at least $N$ have unique fillers
in $\D_NA$, and $\D A$ is the strict limit of the tower $(\D_NA)_N$; it is
also the internal limit. The truncation level is shifted by one: $\D_0A$
is the propositional truncation of $A$ and $\D_1A$ its set truncation, and
in general $\D_{n+1}$ is the candidate for $n$-truncation. The maps of the
tower are not fibrations.
\end{abstract}

%======================================================================
\section{The coskeleton}\label{sec:cosk}

We use the notation of \texttt{localisation.tex}: $\C$ is the category of
finite non-empty posets, $\Delta\subseteq\C$ the linear posets
$[n]=\{0<\dots<n\}$, $\Delta_+\subseteq\Delta$ the strictly monotone maps,
and $i^*$ the restriction to $\Delta$.

\begin{definition}
Let $N\ge-1$ and let $A$ be a type over $\Gamma$. For $\rho\in\Gamma(X)$,
$K_NA(X,\rho)$ is the set of families $c(f)\in A([k],\rho f)$, for all
$k\le N$ and $f:[k]\to X$, such that $c(f)\,g=c(fg)$ for every
$g:[j]\to[k]$ with $j\le N$. Restriction along $X'\to X$ is precomposition.
The map $\eta:A\to K_NA$ is $\eta(a)(f)=af$, and $K_MA\to K_NA$, for
$N\le M$, forgets the values above $N$.
\end{definition}

The family and the condition are coded, so $K_NA$ is a type. On a linear
poset, $K_NA([m],\tau)$ is the set of maps from the $N$-skeleton of
$\Delta[m]$ to $i^*A$ over $\tau$: such a map is determined by its values
on the simplices of dimension at most $N$, and these are natural along all
maps. So $i^*K_NA$ is the $N$-coskeleton of $i^*A$. For $N=-1$ the
families are empty and $K_{-1}A=1$.

\begin{lemma}\label{lem:cosk}
\leavevmode
\begin{enumerate}
\item For $m\le N$, $\eta:A([m],\tau)\to K_NA([m],\tau)$ is a bijection,
  with inverse $c\mapsto c(\id)$.
\item For $m>N$, restriction to the faces is a bijection from
  $K_NA([m],\tau)$ to the compatible families of elements of
  $K_NA([m-1],\tau d_i)$, $0\le i\le m$.
\item A fibration structure on $A$ gives a horn filling operation on
  $i^*K_NA$.
\item On a linear poset, $A$ is the limit of the tower
  $\cdots\to K_{N+1}A\to K_NA\to\cdots\to K_{-1}A$.
\end{enumerate}
\end{lemma}

\begin{proof}
(1) For $m\le N$ the identity of $[m]$ is one of the indices, and
$c(g)=c(\id)\,g$.

(2) A map $[k]\to[m]$ with $k\le N<m$ misses a point of $[m]$, so it
factors through a face $d_i$. If it factors through $d_i$ and $d_j$, it
factors through their intersection, on which compatible faces agree.

(3) Let $h$ be a horn of $i^*K_NA$ on $\Lam^k[m+1]$. If $m\le N$, its faces
read at the identity are elements of $A([m])$, and two of them agree on
their common face, by the naturality of $h$ and (1). The construction of
Lemma~\ref{L-lem:comp-horn} uses only these faces and their agreement: it
reads them on the poset horn and composes in $A$. Its filler $a$ gives the
filler $\eta(a)$. If $m>N$, a map $[j]\to[m+1]$ with $j\le N$ misses at least
two vertices, hence a vertex different from $k$, so it lies in the horn,
and $h$ read at these maps is the filler.

(4) By (1).
\end{proof}

The count in (3) is elementary: a monotone sequence on $[0,j]$ with values
at most $M$ and $j<M$ misses a value at most $M$; applied to the
composite with the degeneracy at $k$, this gives a vertex $y\ne k$
missed by the map.

%======================================================================
\section{The approximations $\D_N$}\label{sec:DN}

\begin{definition}
For a type $A$ over $\Gamma$ and $N\ge-1$, $\D_NA=\D(K_N A)$.
\end{definition}

\begin{proposition}\label{prop:DN}
$\D_NA(X,\rho)$ is in bijection with the set of families $w(f)\in
A([k],\rho f)$, for $k\le N$ and $f:[k]\to X$, with $w(f)\,g=w(fg)$ for
every strictly monotone $g:[m]\to[k]$. The bijection commutes with
restriction.
\end{proposition}

\begin{proof}
Given $u\in\D_NA(X,\rho)$, put $w(f)=u(f)(\id)$; it is natural along strict
maps, by the naturality of $u$ and of $u(f)$. Conversely, given $w$, put
$u(f)(g)=w(fm)\,\sigma$, where $g=m\sigma$ is the image factorisation of
$g:[j]\to[k]$, $m$ strict and $\sigma$ surjective. Two factorisations are
compared by a strict map $\mu$ (Eilenberg--Zilber): for $g$ and $gt$,
$\sigma t=\mu\sigma'$ and $m'=m\mu$, which gives the naturality of $u(f)$
along every $t$; for $g$ and $sg$ with $s$ strict, $\sigma=\mu\sigma'$ and
$m\mu=sm'$, which gives the naturality of $u$ along $s$. Below $N$ the value
$w(fm)\sigma$ is $w(f)\,m\sigma=w(f)g$, so the two maps are inverse.
\end{proof}

So $\D_NA$ is $\D A$ with the indices restricted to $[0],\dots,[N]$, and
$\D A=\D_\infty A$.

\begin{corollary}\label{cor:fib}
If $A$ is fibrant, so is $\D_NA$, naturally in $\Gamma$.
\end{corollary}

\begin{proof}
By Lemma~\ref{lem:cosk}(3) and Theorem~\ref{L-thm:DA} applied to $K_NA$.
Naturality is obtained on the universal family, as in
Theorem~\ref{L-thm:DA}: $K_N$ preserves cartesian maps, as $\D$ does, so
$\D_N(c^*\El)\to\D_N\El$ is cartesian over $c$, and these maps compose.
\end{proof}

\begin{example}\label{ex:low}
\leavevmode
\begin{itemize}
\item $\D_{-1}A=1$.
\item $\D_0A(X,\rho)=\prod_{x\in X}A([0],\rho x)$, since the only strictly
  monotone map $[0]\to[0]$ is the identity. Any two elements $u,v$ at $X$
  are the ends of a line, namely the element of $\D_0A$ at $X\times[1]$
  equal to $u$ on $X\times\{0\}$ and to $v$ on $X\times\{1\}$. So $\D_0A$
  is a proposition, with an element over $\rho$ exactly when every $A([0],\rho x)$
  has one: the propositional truncation of $A$, with no higher inductive
  type.
\item $i^*K_1A$ has the points and the edges of $i^*A$, and every
  compatible boundary of dimension $\ge2$ has a unique filler. Its
  connected components are those of $i^*A$, and its loops are all
  homotopic to the constant ones: $\D_1A$ is the candidate set truncation.
\end{itemize}
\end{example}

\begin{remark}
The fibre of $\D_NA$ at $X=[0]$ is not the truncation of the fibre of $A$.
It is the set of maps from the $N$-skeleton of the semi-simplicial set
with one simplex in each dimension, which has the homotopy type of a
sphere when $N$ is odd. This does not matter, because the homotopy type
of a $\D$-modal type is read on its restriction to $\Delta$.
\end{remark}

%======================================================================
\section{Truncation level and limit}\label{sec:trunc}

\begin{proposition}\label{prop:cosk}
For $m>N$, restriction to the faces is a bijection from $\D_NA([m],\rho)$ to
the compatible families of elements of $\D_NA([m-1],\rho d_i)$. So
$i^*\D_NA$ has unique fillers for all spheres of dimension at least $N$.
\end{proposition}

\begin{proof}
By Proposition~\ref{prop:DN}, as in Lemma~\ref{lem:cosk}(2): a map
$[k]\to[m]$ with $k\le N<m$ misses a point, so it factors through a face.
\end{proof}

In a Kan complex with this property, the homotopy groups $\pi_k$ vanish
for $k\ge N$. A class in $\pi_k$ is a $k$-simplex with boundary at the base
point, and the sphere of dimension $k\ge N$ formed by this simplex and
degenerate ones has a filler, which is a homotopy to the base point.
Classically, $i^*A\to\cosk_Ni^*A$ is moreover a bijection on $\pi_k$ for
$k<N$, since these use simplices of dimension at most $k+1\le N$ only. So
$\D_{n+1}A$ is the candidate $n$-truncation of $\D A$.

\begin{proposition}\label{prop:lim}
$\D A$ is the limit of the tower $\cdots\to\D_{N+1}A\to\D_NA\to
\cdots\to\D_{-1}A=1$.
\end{proposition}

\begin{proof}
By Proposition~\ref{prop:DN}, a compatible sequence of elements of the
$\D_NA(X,\rho)$ is a family indexed by all the maps $[k]\to X$, natural
along the strictly monotone maps, that is, an element of $\D A(X,\rho)$.

\end{proof}

\begin{remark}\label{rem:notfib}
The maps of the tower are not fibrations. A line of $\D_0A$ at $[0]$ is an
element at $[1]$, that is, two points $a,b$ of $A([0])$. A lift to $\D_1A$
starting from a point of $\D_1A$ over $a$ gives an element of $\D_1A([1])$,
and in particular an edge of $i^*A$ from $a$ to $b$. If $a$ and $b$ lie in
different components, there is none. The limit is therefore a strict
limit, not the limit of a tower of fibrations. The same holds for the
coskeletal tower of a Kan complex, which converges strictly although its
maps are not Kan fibrations. Internally nothing is lost: every map is
equivalent to the projection of the family of its fibres, so the internal
tower $\D_{N+1}A\to\D_NA$ is a tower of fibrations up to equivalence. What
does not hold automatically is that the strict limit computes the internal limit.
\end{remark}

%======================================================================
\section{Convergence}\label{sec:conv}

Write $C=i^*A$ and $C_N=i^*K_NA$, the $N$-coskeleton of $C$, so that
$\D_NA=\D(K_NA)$. Each $C_N$ has a horn filling operation
(Lemma~\ref{lem:cosk}(3)).
In dimensions $\le N$ the simplices of $C_N$ are those of $C$, and in higher
dimensions they are determined by their boundaries. For a poset $P$ we
identify $P$ with its nerve, so that $\Delta[m]\times\Delta[1]^k$ is the
poset $[m]\times[1]^k$. A map $P\to C_N$ over $\rho$ is an element of
$K_NA(P,\rho)$, since both are the maps from the $N$-skeleton of $P$ to $C$.

\begin{lemma}\label{lem:box}
Let $T$ be a simplicial set with a horn filling operation, $P$ a finite
poset, $S\subseteq P$ a simplicial subset and $e\in\{0,1\}$. Every map
$P\times\{e\}\cup S\times\Delta[1]\to T$ extends to $P\times\Delta[1]$.
In particular every open box of $\Delta[m]\times\Delta[1]^k$, $k\ge1$, has a
filler in $T$.
\end{lemma}

\begin{proof}
\emph{Prisms over a simplex.} A problem at $n$ is a map
$b:\Delta[n]\times\{e\}\cup Q\times\Delta[1]\to T$ with $Q\subseteq\Delta[n]$; its
restriction along $d_i$ is a problem at $n-1$. By recursion on $n$ we choose a
filler $F(b)$ with $F(b)\,d_i=F(b\,d_i)$. If $Q$ contains the identity,
$F(b)=b$. Otherwise the fillers $F(b\,d_i)$ agree on the common faces, since
$F(b\,d_i)\,d_j=F(b\,d_id_j)$, and $\partial\Delta[n]\times\Delta[1]\cup
\Delta[n]\times\{e\}$ is given. The non-degenerate $(n+1)$-simplices of
$\Delta[n]\times\Delta[1]$ are the shuffles
\[
P_k=\bigl((0,0),\dots,(k,0),(k,1),\dots,(n,1)\bigr),\qquad 0\le k\le n.
\]
The face $d_k$ of $P_k$ is the face $d_k$ of $P_{k-1}$, the face $d_{k+1}$ of
$P_k$ is the face $d_{k+1}$ of $P_{k+1}$, and the other faces of $P_k$ lie in
$\partial\Delta[n]\times\Delta[1]$, apart from $d_{n+1}P_n=\Delta[n]\times\{0\}$
and $d_0P_0=\Delta[n]\times\{1\}$. For $e=0$ fill $P_n,P_{n-1},\dots,P_0$ in
this order: when $P_k$ is filled, all its faces but $d_k$ are known, so $P_k$
is the filler of a horn on $\Lam^k[n+1]$. For $e=1$ fill $P_0,\dots,P_n$, each
on $\Lam^{k+1}[n+1]$.

\emph{Prisms over a poset.} For an injective $c:[n]\to P$ let $b_c$ be the
problem at $n$ read along $c$, with $Q=c^*S$. Restriction along a face map
gives $b_c\,d_i=b_{cd_i}$, so $F(b_c)\,\mu=F(b_{c\mu})$ for every injective
$\mu$, since these are composites of face maps. A simplex $g:[W]\to P$ factors uniquely as
$g=c\sigma$ with $c$ injective and $\sigma$ surjective (Eilenberg--Zilber), and
we put $F(g,j)=F(b_c)(\sigma,j)$. If $k:[W']\to[W]$, the factorisation of
$gk$ is $(c\mu)\sigma'$ with $\mu$ injective and $\mu\sigma'=\sigma k$, so
$F(gk,jk)=F(b_c)(\mu\sigma',jk)=F(g,j)\,k$. If $g\in S$ then $c\in S$, since
$c=g\delta$ for a section $\delta$ of $\sigma$; the identity is then in
$c^*S$, and $F(g,j)=b_c(\sigma,j)$ is the given value. The same holds at
$j=e$.

\emph{Boxes.} An open box of $\Delta[m]\times\Delta[1]^k$ is the union of all the
codimension-one faces except one face $x_l=1-e$ of the cube. Permuting the cube
coordinates so that $l=k$, it is $P\times\{e\}\cup\partial P\times\Delta[1]$,
where $P=[m]\times[1]^{k-1}$ and $\partial P$ is the union of the faces of $P$.
\end{proof}

The fillers are uniform along the face maps only, so no excluded middle is
used. The prism over $P$ is filled simplex by simplex, without ordering
the simplices of $P$. The shuffle order is used only inside each simplex.
Let
\[
L=\Sigma\,(u:\textstyle\prod_N\D_NA)\,\prod_N\Path_{\D_NA}(r\,u_{N+1},u_N)
\]
be the internal limit, and $\varphi:\D A\to L$ the map
$u\mapsto((r\,u)_N,\text{constant lines})$.

\begin{theorem}\label{thm:conv}
If $A$ is fibrant, $\varphi$ is a homotopy equivalence.
\end{theorem}

\begin{proof}[Construction]
Let $\ell=(u_N,p_N)\in L(X)$. For $f:[m]\to X$ put $\sigma_N=u_N(f)$, and
let $p_N(f)$ be the prism $\Delta[m]\times\Delta[1]\to C_N$ obtained by
reading $p_N$ along $f\times[1]$ and extending by Eilenberg--Zilber, as in
Theorem~\ref{L-thm:DA}. We define, by recursion on $m$, an
$m$-simplex $\psi(f)$ of $C$, and for every $N\ge m$:
\begin{itemize}
\item a prism $H^N(f):\Delta[m]\times\Delta[1]\to C_N$ from $\psi(f)$ to
  $\sigma_N$;
\item a square $\Theta_N(f):\Delta[m]\times\Delta[1]^2\to C_N$ with edges
  $H^{N+1}(f)$ and $H^N(f)$, $p_N(f)$, and the constant line at $\psi(f)$.
\end{itemize}
Each of these restricts along the face maps to the same datum at $fd_i$.
For $N<m$ there is nothing to choose, since prisms and squares in $C_N$ of
an $m$-simplex with $m>N$ are determined by their faces. All the open boxes
below are filled by Lemma~\ref{lem:box}. Each lies in dimensions $\le m+1$, so in
step 3, where $m+1\le N-1$, the box given in $C_{N-1}$ is a box of $C_N$, and
the image in $C_{N-1}$ of the filled square is $\Theta_{N-1}(f)$. The boxes of
steps 1 and 3 have dimension at most $N$, so they are boxes of $C$ itself, and
only step 2 uses $C_m$. Steps 1 and 3 alone define $\psi(\ell)$; step 2 is
needed only for the homotopy.
\begin{enumerate}
\item $H^{m+1}(f)$ is the filler of the open box with lid $\sigma_{m+1}$
  and sides $H^{m+1}(fd_i)$. Its other end is $\psi(f)$. The box lies in
  dimensions $\le m+1$, where $C_{m+1}$ is $C$, so $\psi(f)$ is a simplex
  of $C$.
\item $H^m(f)$ and $\Theta_m(f)$ fill, in $C_m$, the open box of the square
  formed by $H^{m+1}(f)$, $p_m(f)$, the constant line, and the sides
  $\Theta_m(fd_i)$.
\item For $N\ge m+2$, $H^N(f)$ and $\Theta_{N-1}(f)$ fill, in $C_N$, the
  open box formed by $H^{N-1}(f)$, $p_{N-1}(f)$, the constant line, and the
  sides $\Theta_{N-1}(fd_i)$.
\end{enumerate}
The data of $fd_i$ used in each step have already been built. For $fd_i$,
of dimension $m-1$, they come from its steps 1 and 3 with $N=m$ and
$N=m+1$. The edges shared by a side and the rest of the box agree because
they are themselves sides of earlier boxes. So $\psi(\ell)=(f\mapsto\psi(f))$
is natural along the face maps, and is an element of $\D A(X)$. Each
$(H^N,\Theta_N)$ is natural along the face maps, so it defines a line of
$\D_NA$ and a square, read at the maps $(f,h):[k]\to X\times[1]$ through
$\langle\id,h\rangle$, as in Lemma~\ref{L-lem:D-htpy}. These form a path
$\varphi\psi(\ell)\simeq\ell$ in $L$. Everything at $f$ is computed from
data along $f$, so the construction is natural in $X$ and in $\Gamma$.

For $\psi\varphi\simeq\id$, let $u\in\D A(X)$ and run the construction on
$\varphi(u)$. By the same recursion, using open boxes one dimension
higher, build for each $f:[m]\to X$:
\begin{itemize}
\item a prism $G(f)$ in $C$ from $\psi\varphi(u)(f)$ to $u(f)$, with sides
  $G(fd_i)$;
\item for $N\ge m+1$, squares $\Xi^N(f)$ in $C_N$ between $H^N(f)$ and
  $G(f)$, with constant top and bottom.
\end{itemize}
$G(f)$ and $\Xi^{m+1}(f)$ fill the box of $H^{m+1}(f)$ with sides
$\Xi^{m+1}(fd_i)$. For $N\ge m+2$, $\Xi^N(f)$ fills, in $C_N$, the cube
formed by $\Theta_{N-1}(f)$, $\Xi^{N-1}(f)$, constant faces, and the sides
$\Xi^N(fd_i)$. Then $G$ is a line of $\D A$ from $\psi\varphi(u)$ to $u$.
\end{proof}

\begin{remark}
Two features make the construction work.
\begin{itemize}
\item The value of $\psi(\ell)$ in dimension $m$ is taken at stage $m+1$,
  never at a lower stage. In particular $u_0$ is never used, so the
  obstruction of Remark~\ref{rem:notfib} does not arise.
\item The squares $\Theta_N$ are part of the data and are built by box
  filling. They are not obtained by concatenating the $p_N$. A
  concatenation would not commute strictly with the faces: the face of
  $(H\cdot p^{-1})\cdot p$ is not $H$.
\end{itemize}
The connectedness of the maps of the tower is witnessed by fillers of
$C$ and of the $C_N$. So no choice principle is used, whereas the internal
argument from towers of connected maps needs dependent choice.
\end{remark}

%======================================================================
\section{The propositional truncation}\label{sec:prop}

A type $B$ is a proposition if any two elements $w,w'$ of $B$ over the same
point have a path, natural in the stage: an element of $B$ at $X\times[1]$
equal to $w$ on $X\times\{0\}$ and to $w'$ on $X\times\{1\}$. This is a
section of $\Pi(b\,b':B)\,\Path_B(b,b')$, read in the total presheaf.

\begin{theorem}\label{thm:prop}
Let $A$ be fibrant. Then $\D_0A$ is a fibrant $\D$-modal proposition, and
for every fibrant $\D$-modal proposition $B$ and map $g:A\to B$ there is a
map $\D_0A\to B$. So $A\to\D_0A$ is the propositional truncation in the
$\D$-modal types.
\end{theorem}

\begin{proof}
$\D_0A$ is fibrant by Corollary~\ref{cor:fib}, and modal since $\D C$ is
modal whenever $i^*C$ has a horn filling operation, which holds for
$C=K_0A$ (Lemma~\ref{lem:cosk}(3)). It is a proposition by
Example~\ref{ex:low}: the path between $u$ and $v$ is $u$ on $X\times\{0\}$
and $v$ on $X\times\{1\}$.

A fibrant proposition $B$ has sphere fillers, uniform along the face maps.
For $w\in\D_0B([k],\rho)$, that is $k+1$ vertices of $B$ over a simplex
$\rho$ of the context, define $\mathrm{fl}(w)\in B([k],\rho)$ with faces
$\mathrm{fl}(w\,d_i)$, by recursion on $k$. For $k=0$ it is the vertex. For
$k+1$ the faces $s_i=\mathrm{fl}(w\,d_i)$ form a sphere. A filler $c$ of the
horn $(s_i)_{i\ge1}$ (Lemma~\ref{L-lem:comp-horn}) lies over $\rho$; the box
over $[k+1]$ with lid $c$ at $1$ and, on the boundary, the paths of $B$
from the sphere to $c$ is filled by Lemma~\ref{lem:box}, and
$\mathrm{fl}(w)$ is its other end. Then $f\mapsto\mathrm{fl}(w f)$ is a map
$\D_0B\to\D B$, natural since $\mathrm{fl}$ depends only on $w f$. The map
$\D_0A\to B$ is $p_B\circ\mathrm{fl}\circ\D_0 g$, with $p_B$ an inverse of
$\eta_B$. The homotopies needed for the universal property are free,
since $B$ is a proposition.
\end{proof}

For $n\ge0$ the same pattern needs sphere fillers in an $n$-truncated type
for spheres of dimension $\ge n+1$, which requires relating iterated path
types to spheres (Section~\ref{sec:remains}, item 1). The alternative
induction through $\Path_{\D_{n+2}A}(\eta a,\eta b)\simeq\D_{n+1}\Path_{\D
A}(\eta a,\eta b)$ needs, besides this comparison (a homotopy equivalence,
not an isomorphism: the $(n+2)$-simplices of $\D_{n+2}A$ over a
non-degenerate $(n+2)$-simplex of $X$ have no counterpart), the internal
lemmas on connected maps, for which the model has no internal language.

%======================================================================
\section{What remains}\label{sec:remains}

The statement \emph{every $\D$-modal type is the limit of its Postnikov
tower} needs the following points beyond Sections~\ref{sec:cosk}--\ref{sec:conv}.
\begin{enumerate}
\item The comparison between the internal truncation level (iterated
  identity types are propositions) and the unique filling of spheres of
  Proposition~\ref{prop:cosk}, for fibrant types of the poset model.
\item The universal property of $A\to\D_{n+1}A$ among the $\D$-modal
  $n$-truncated types $B$. A natural route is to show that $\D B\to\D_{n+1}B$ is a
  homotopy equivalence for such $B$, since spheres of dimension $>n$ bound
  in $B$, and then to argue as in Proposition~\ref{L-prop:univ}. That
  argument also needs $\eta_{\D_NA}\simeq\D_N\eta_A$, and the homotopy of
  Lemma~\ref{L-lem:eta-path} does not restrict to $\D_N$: it reads $u$ at
  dimensions up to $|B_0|+|A_1|-1$, which may exceed $N$.
\item Whitehead's principle in the form: if $\alpha:A\to B$ is a map of
  fibrant types and every $\D_N\alpha$ is a homotopy equivalence, then so is
  $\D\alpha$. With Theorem~\ref{thm:conv} it reduces to the invariance of the
  internal limit under a levelwise equivalence of towers. That invariance
  is a statement of type theory with $\Pi$, $\Sigma$ and $\Path$, and needs
  no property of the model.
\end{enumerate}

%======================================================================
\appendix
\section{Correspondence with the formalisation}

\noindent Layer L9 of \texttt{seg/core}, over the poset instance.

\medskip
\begin{tabular}{p{4.3cm}p{9.6cm}}
statement & files\\\hline
$K_NA$, $\eta$, the restrictions & CCosk (KT, Res)\\
Lemma~\ref{lem:cosk}(3) & CCoskHorn (HK; the count: pigN, InHorn); CLocHorn now reads a horn through its faces (FD, HSF, HPF, FaceFill), and HornOfFib is the instance for $i^*A$\\
Definition of $\D_N$, Corollary~\ref{cor:fib} & CLocDN (DN.fib = Theorem~\ref{L-thm:DA} for $K_NA$); naturality: CCoskCart, CLocDNUniv\\
Proposition~\ref{prop:DN} & CLocDNC (to, from, to-from, from-to, to-nat); CLocEZ.sN-le, r-le\\
Proposition~\ref{prop:lim} & CLocDN (Lim: u, RN.res-u, uniq)\\
Lemma~\ref{lem:box} & CBoxFill (BF: the problem at $c$, RS, the filled prisms fl, the map Fv with Fp, Fn, Fe, Fs); the factorisation: CBoxEZ (EZX, EZXK, EZXI; heights HLin, HProd); faces suffice: CFacesGen (CLocD.FacesGen is its instance); squares: CBoxSq\\
Theorem~\ref{thm:conv}, the construction of $\psi$ & reading $u_N$, $p_N$ along $f$: CPostMap (DR, DRI, DRZ, DRP, DRId); sides of the boxes: CPostBdry (BG, BS; CFaceTwo, shared with CLocHorn); the open box of a square over a simplex, for any type with face-uniform fillers: CPostSq (SqE, SqS, SqZ); step 1: CPostStep1 (P1, S1); step 3: CPostStep3 (E3; S3, instance of SqS); dimension 0: CPostStep0; the recursion: CPostRec (Rec, SRec over the stages, ST from $m$ to $m+1$, FFS the agreement of faces, Drec, psi $\in\D A(X)$)\\
Theorem~\ref{thm:prop} & IsPropE, D0P (D$_0$A a proposition), D0M (fibrant, modal; CLocUniv.dModalH, CLocEq.FibD): CTruncProp; SF (sphere filling fl, the map sN : $\D_0B\to\D B$), KMap ($K_N$ on maps), UP0.k ($\D_0A\to B$): CTruncFill\\
Theorem~\ref{thm:conv}, $\varphi\psi\simeq\id$ & $\eta$ and the value at the identity (Lemma~\ref{lem:cosk}(1)): CCosk (etaE, ev, eta-ev, ev-eta); reading $p_N$ along $f\times[1]$ in $K_NA$: CPostK (DRK, DRKI, DRKE; along strict chains DRS, DRKS); step 2: CPostStep2 (E2; S2, S2z, instances of SqS, SqZ for $B=K_mA$ with the fillers of CCoskHorn); faces suffice below a bound: CFacesGen.SGB (SG is its instance); the lines: CPostLine (St2S, St2Z: step 2 at every simplex; the family val, natural along faces WF, along strict maps FG; DCN.line $\in\D_NA(X\times[1])$; end0, end1); the squares: CPostSquare (val2, DCS.sq $\in\D_NA((X\times[1])\times[1])$; edges sq-s0, sq-s1, sq-t0, sq-t1)
\end{tabular}

\medskip
Not formalised: Lemma~\ref{lem:cosk}(2) and Proposition~\ref{prop:cosk},
Example~\ref{ex:low} (except $\D_0A$ a proposition), Section~\ref{sec:remains}; of Theorem~\ref{thm:conv},
the homotopy $\psi\varphi\simeq\id$ and the naturality of $\psi$ and of the
homotopy in $X$. For $\varphi\psi\simeq\id$ the lines and squares are
formalised with all their edges, at a fixed $(X,\ell)$; the internal limit $L$
is not defined as a type, and the path in $L$ is given by these components.

\end{document}
